\phantomsection
\addcontentsline{toc}{chapter}{DANH MỤC CHỮ VIẾT TẮT, THUẬT NGỮ CHUYÊN MÔN}

\begin{center}
{\large\bfseries{DANH MỤC CHỮ VIẾT TẮT, THUẬT NGỮ CHUYÊN MÔN}}
\end{center}

\vspace{1em}

\noindent\textbf{1. Chữ viết tắt}

{
\small
\renewcommand{\arraystretch}{1.2}
\begin{longtable}{L{2.6cm}L{4.8cm}L{6.3cm}}
\toprule
\textbf{Ký hiệu} & \textbf{Cụm từ tiếng Anh} & \textbf{Giải thích} \\
\midrule
\endfirsthead
\toprule
\textbf{Ký hiệu} & \textbf{Cụm từ tiếng Anh} & \textbf{Giải thích} \\
\midrule
\endhead
SBD & Shot Boundary Detection & Phát hiện ranh giới cảnh quay trong video. \\
SHOT & Short Video Shot Boundary Detection Dataset & Bộ dữ liệu chuyên biệt cho bài toán SBD trên video ngắn. \\
NAS & Neural Architecture Search & Tìm kiếm kiến trúc mạng nơ-ron; kỹ thuật tự động tìm kiến trúc mạng tối ưu. \\
SPOS & Single-Path One-Shot & Chiến lược NAS một lần theo đường đơn với cơ chế lấy mẫu đường đi đơn. \\
CNN & Convolutional Neural Network & Mạng nơ-ron tích chập dùng để học đặc trưng hình ảnh. \\
3D ConvNet & 3D Convolutional Network & Mạng tích chập ba chiều khai thác đồng thời thông tin không gian và thời gian. \\
DDCNN & Dilated Deep Convolutional Neural Network & Khối tích chập giãn nở dùng trong TransNetV2 và AutoShot. \\
TP & True Positive & Số chuyển cảnh được phát hiện đúng. \\
FP & False Positive & Dương tính giả; số báo động giả. \\
FN & False Negative & Số chuyển cảnh có thật nhưng bị bỏ sót. \\
BCE & Binary Cross-Entropy & Hàm mất mát entropy chéo nhị phân cho bài toán phân loại hai lớp. \\
KD & Knowledge Distillation & Chưng cất tri thức từ mô hình giáo viên sang mô hình học sinh. \\
NLL & Negative Log-Likelihood & Âm log-hợp lý; hàm mục tiêu thường dùng khi hiệu chỉnh xác suất bằng hiệu chỉnh nhiệt độ. \\
EMA & Exponential Moving Average & Trung bình trượt lũy thừa trọng số mô hình qua các bước huấn luyện, giúp làm trơn và ổn định trọng số cuối. \\
CV & Cross-Validation & Kiểm định chéo; khóa luận dùng kiểm định chéo 5 phần để kiểm tra xu hướng kết quả. \\
F1 & F1-score & Chỉ số điều hòa cân bằng giữa độ chính xác dương (Precision) và độ bao phủ (Recall). \\
\bottomrule
\end{longtable}
}

\noindent\textbf{2. Thuật ngữ chuyên môn}

{
\small
\renewcommand{\arraystretch}{1.2}
\begin{longtable}{L{4.2cm}L{8.8cm}}
\toprule
\textbf{Thuật ngữ tiếng Anh} & \textbf{Giải thích bằng tiếng Việt} \\
\midrule
\endfirsthead
\toprule
\textbf{Thuật ngữ tiếng Anh} & \textbf{Giải thích bằng tiếng Việt} \\
\midrule
\endhead
Shot & \textbf{Cảnh quay:} Đoạn cảnh quay liên tục giữa hai ranh giới chuyển cảnh. \\
Scene & \textbf{Cảnh:} Tập hợp các cảnh quay có cùng nội dung hoặc không gian, thời gian. \\
Frame Similarity & \textbf{Độ tương đồng khung hình:} Mức độ tương đồng giữa các khung hình kế cận để phát hiện chuyển cảnh. \\
Color Histogram & \textbf{Biểu đồ màu:} Biểu đồ thống kê biểu diễn phân bố màu sắc của khung hình. \\
Factorized Convolution & \textbf{Tích chập phân tách:} Tách tích chập ba chiều thành các bước xử lý không gian và thời gian độc lập. \\
Skip Connection & \textbf{Kết nối tắt:} Kết nối bỏ qua một hoặc nhiều tầng mạng giúp ổn định lan truyền gradient. \\
Focal Loss & \textbf{Hàm mất mát Focal:} Hàm mất mát tập trung vào các mẫu khó phân loại và giảm trọng số mẫu dễ. \\
Many-hot Label & \textbf{Nhãn nhiều điểm dương:} Nhãn mở rộng vùng biên quanh ranh giới để tăng tín hiệu giám sát. \\
Backbone & \textbf{Mạng xương sống:} Thành phần trích xuất đặc trưng cơ bản từ đầu vào. \\
Temporal Modeling & \textbf{Mô hình hóa thời gian:} Mô hình hóa các biến đổi của đặc trưng hình ảnh theo thời gian. \\
Soft Label & \textbf{Nhãn mềm:} Nhãn dạng xác suất mang giá trị liên tục từ 0 đến 1 thay vì nhãn cứng. \\
Temperature Scaling & \textbf{Hiệu chỉnh nhiệt độ:} Hiệu chỉnh xác suất đầu ra bằng cách chia logit cho một hệ số nhiệt độ. \\
Calibration & \textbf{Hiệu chỉnh độ tin cậy:} Quá trình làm cho xác suất đầu ra khớp với độ tin cậy thực tế của dự đoán. \\
Logit & \textbf{Logit:} Giá trị đầu ra thô của mô hình trước khi đi qua hàm kích hoạt xác suất. \\
Gaussian Smoothing & \textbf{Làm mượt Gaussian:} Phép lọc làm mượt chuỗi xác suất theo thời gian để giảm thiểu đỉnh nhiễu. \\
Decision Threshold & \textbf{Ngưỡng quyết định:} Ngưỡng xác suất để đưa ra quyết định nhị phân về sự tồn tại của ranh giới. \\
Threshold Sweep & \textbf{Quét ngưỡng:} Quá trình dò tìm ngưỡng quyết định tối ưu dựa trên các chỉ số đánh giá. \\
SuperNet & \textbf{Siêu mạng:} Kiến trúc mạng tích hợp chứa toàn bộ các mạng con ứng viên trong NAS. \\
Bayesian Optimization & \textbf{Tối ưu hóa Bayesian:} Phương pháp tối ưu hóa hộp đen bằng cách xây dựng một mô hình xác suất. \\
Weight Grafting & \textbf{Ghép trọng số:} Phương pháp ghép nối các lớp trọng số tối ưu từ nhiều mô hình khác nhau. \\
DARTS & \textbf{Tìm kiếm kiến trúc khả vi:} Phương pháp tìm kiếm kiến trúc mạng nơ-ron bằng tối ưu hóa liên tục khả vi. \\
Domain Adaptation & \textbf{Thích ứng miền:} Phương pháp giúp mô hình hoạt động ổn định khi phân phối dữ liệu thay đổi miền. \\
Dropout & \textbf{Loại bỏ ngẫu nhiên:} Phương pháp chính quy hóa tắt ngẫu nhiên các nút mạng khi huấn luyện để tránh quá khớp. \\
Baseline & \textbf{Mô hình cơ sở:} Mô hình đối chứng được sử dụng để so sánh và đánh giá hiệu năng. \\
Checkpoint & \textbf{Tệp trọng số mô hình:} Tệp lưu trữ các giá trị trọng số của mô hình trong quá trình huấn luyện. \\
Weight Decay & \textbf{Suy giảm trọng số:} Kỹ thuật chính quy hóa phạt các trọng số lớn để hạn chế quá khớp. \\
Ablation Study & \textbf{Thực nghiệm phân rã:} Thực nghiệm cô lập hoặc loại bỏ thành phần để đo lường đóng góp của từng thành phần. \\
Grid Search & \textbf{Tìm kiếm theo lưới:} Phương pháp quét qua một tập hợp siêu tham số định sẵn để tìm cấu hình tốt nhất. \\
Evaluation Protocol & \textbf{Giao thức đánh giá:} Bộ quy ước chuẩn về cách tính toán, đo đạc và báo cáo hiệu năng mô hình. \\
Heatmap Label & \textbf{Nhãn dạng bản đồ nhiệt:} Nhãn xác suất phân phối Gaussian quanh vùng biên để hỗ trợ học chuyển cảnh dần. \\
\bottomrule
\end{longtable}
}
